%% toc: P2 13 13.1,13.2,13.3
\section{Derivation maps between the Hilbert, natural-deduction and sequent calculi}
\label{sec:p02-w02b-p02-010-translations}

The three classical propositional calculi of this part prove the same
consequences, but they do so with different objects: the Hilbert calculus
\(\mathsf H_{\mathsf{CPL}}\) of \cref{def:p02-hilbert-derivation} produces
finite sequences of \(\{\neg,\to\}\)-formulas, the natural-deduction calculus
\(\mathsf{ND}_{\mathsf{CPL}}\) produces trees with discharged assumptions,
and \(\mathsf{LK}_{\mathsf{CPL}}\) produces trees of sequents.  This section
gives explicit constructions that turn a derivation in one calculus into a
derivation in another, in every direction among the three, and proves that
each construction is correct.  The existence statements obtained from them
are the Hilbert legs of \cref{thm:p02-calculi-derivability-equivalence}; the
constructions themselves say more, because they fix which derivation is
produced.

\emph{Inputs.}  A Hilbert derivation from a set of \(\{\neg,\to\}\)-formulas;
or an \(\mathsf{ND}_{\mathsf{CPL}}\) derivation over the full language with
primitives \(\bottom,\neg,\land,\lor,\to\); or an \(\mathsf{LK}_{\mathsf{CPL}}\)
derivation of a sequent of full-language formulas.  \emph{Conclusions.}  A
derivation in each of the other two calculi of the same formula, when it
lies in the \(\{\neg,\to\}\)-fragment, or otherwise of its image under the
translation \(\operatorname{tr}_{\mathsf H}\) of
\cref{def:p02-hilbert-translation}, from the correspondingly translated
assumptions.  The constructions from natural deduction to sequents and back
are those of \cref{thm:p07-lk-nd-equivalence}; this section supplies the
four constructions that involve the Hilbert calculus and checks that the
sequent constructions stay inside the propositional calculi.

\subsection{Setting}
\label{sec:p02-w02b-p02-010-setting}

\begin{hbconvention}{def:p02-w02b-p02-010-setting}{Calculi, open assumptions and the two fixed formulas}
\(\mathsf{Fm}_{\mathsf{CPL}}\) is the set of propositional formulas with
primitives \(\bottom,\neg,\land,\lor,\to\), and
\(\mathsf{Fm}_{\{\neg,\to\}}\subseteq\mathsf{Fm}_{\mathsf{CPL}}\) is the set
of formulas built from propositional variables by \(\neg\) and \(\to\) alone.
\begin{enumerate}
  \item \(\Gamma\derives_{\mathsf H_{\mathsf{CPL}}}\theta\) is the relation of
    \cref{def:p02-hilbert-derivation}, for
    \(\Gamma\cup\{\theta\}\subseteq\mathsf{Fm}_{\{\neg,\to\}}\), with axiom
    schemas A1--A3 and the rule MP of
    \texttt{sheet:\allowbreak hilbert-cpl}.  An
    \emph{\(\mathsf H\)-theorem} is a formula derivable from \(\varnothing\).
  \item An \(\mathsf{ND}_{\mathsf{CPL}}\) derivation is a tree of the
    inherited derivation definition, with the rules of
    \texttt{sheet:\allowbreak nd-cpl}; every discharge is
    nonvacuous and applies to all open occurrences of its formula.  For a
    derivation \(D\), \(\operatorname{oa}(D)\) is the set of formulas that
    have at least one open assumption occurrence in \(D\).  For an arbitrary
    \(\Gamma\subseteq\mathsf{Fm}_{\mathsf{CPL}}\),
    \(\Gamma\derives_{\mathsf{ND}_{\mathsf{CPL}}}\varphi\) holds iff some
    derivation \(D\) has root \(\varphi\) and
    \(\operatorname{oa}(D)\subseteq\Gamma\).
  \item An \(\mathsf{LK}_{\mathsf{CPL}}\) derivation is a finite tree of
    sequents of finite multisets of full-language formulas, built from the
    contexted initial sequents by the structural, logical and cut rules of
    the inherited classical sequent calculus, which are the propositional
    rows of \texttt{sheet:\allowbreak lk}.  A finite set written
    as an antecedent is the multiset with multiplicity one on each member.
  \item Fix the propositional variable \(p_0\) of
    \cref{def:p02-hilbert-translation}, and put
    \[
      \top_{\mathsf H}\defeq p_0\to p_0,
      \qquad
      \bot_{\mathsf H}\defeq\neg\top_{\mathsf H}
      =\operatorname{tr}_{\mathsf H}(\bottom).
    \]
    These are names for two fixed formulas of the fragment, not new
    connectives.  By \cref{thm:p02-w02b-p02-010-hilbert-theorems}(1),
    \(\top_{\mathsf H}\) is an \(\mathsf H\)-theorem.
\end{enumerate}
Throughout, \(\varphi,\psi,\chi,\theta\) and \(A,B,C\) range over
\(\mathsf{Fm}_{\{\neg,\to\}}\) when they occur in Hilbert statements, and
\(\alpha,\beta,\gamma\) range over \(\mathsf{Fm}_{\mathsf{CPL}}\).
\end{hbconvention}

The fragment is literally a subset of the full language, and
\(\operatorname{tr}_{\mathsf H}\) is the identity on it: a structural
induction on \(\alpha\in\mathsf{Fm}_{\{\neg,\to\}}\) uses only the clauses for
variables, \(\neg\) and \(\to\), each of which returns the same constructor
applied to the translated parts.  Hence the Hilbert-to-other directions need
no formula translation, while the directions into the Hilbert calculus must
pass through \(\operatorname{tr}_{\mathsf H}\).

\subsection{Hilbert axiom schemas, modus ponens and the deduction theorem}
\label{sec:p02-w02b-p02-010-hilbert}

The calculus has the three schemas and one rule displayed on
\texttt{sheet:\allowbreak hilbert-cpl}: A1
\(\varphi\to(\psi\to\varphi)\), A2
\((\varphi\to(\psi\to\chi))\to((\varphi\to\psi)\to(\varphi\to\chi))\), A3
\((\neg\psi\to\neg\varphi)\to((\neg\psi\to\varphi)\to\psi)\), and MP from
\(\varphi\to\psi\) and \(\varphi\) to \(\psi\).  Every instance of A1--A3 is
an axiom; a member of \(\Gamma\) may be written at any line; MP is the only
rule, and it has no side condition.  The deduction theorem,
\cref{thm:p02-hilbert-deduction}, is proved for arbitrary \(\Gamma\) in
\cref{annex:p2.calc.2} from A1, A2 and MP alone.  The following lemmas are
the Hilbert facts on which the maps below depend.

\begin{hblemma}{thm:p02-w02b-p02-010-hilbert-structural}{proved=annex:p02-w02b-p02-010-hilbert-structural}{Structural properties of Hilbert derivability}
Let \(\Gamma,\Delta\subseteq\mathsf{Fm}_{\{\neg,\to\}}\) and
\(\theta\in\mathsf{Fm}_{\{\neg,\to\}}\).
\begin{enumerate}
  \item (Monotonicity.) If \(\Gamma\derives_{\mathsf H_{\mathsf{CPL}}}\theta\)
    and \(\Gamma\subseteq\Delta\), then
    \(\Delta\derives_{\mathsf H_{\mathsf{CPL}}}\theta\).
  \item (Composition.) If \(\Gamma\derives_{\mathsf H_{\mathsf{CPL}}}\delta\)
    for every \(\delta\in\Delta\) and
    \(\Gamma\cup\Delta\derives_{\mathsf H_{\mathsf{CPL}}}\theta\), then
    \(\Gamma\derives_{\mathsf H_{\mathsf{CPL}}}\theta\).
  \item (Finiteness.) If \(\Gamma\derives_{\mathsf H_{\mathsf{CPL}}}\theta\),
    then \(\Gamma_0\derives_{\mathsf H_{\mathsf{CPL}}}\theta\) for the finite
    set \(\Gamma_0\subseteq\Gamma\) of members of \(\Gamma\) that occur as
    lines of a given derivation.
\end{enumerate}
\end{hblemma}

\begin{hblemma}{thm:p02-w02b-p02-010-hilbert-theorems}{proved=annex:p02-w02b-p02-010-hilbert-theorems}{Derived Hilbert theorems}
For all \(\varphi,\psi,\chi\in\mathsf{Fm}_{\{\neg,\to\}}\), each of the
following has an explicit derivation in \(\mathsf H_{\mathsf{CPL}}\) from the
indicated premises (none, for items 1 and 3--8):
\begin{enumerate}
  \item \(\varphi\to\varphi\);
  \item \(\varphi\to\chi\) from \(\{\varphi\to\psi,\ \psi\to\chi\}\);
  \item \(\neg\neg\varphi\to\varphi\);
  \item \(\varphi\to\neg\neg\varphi\);
  \item \(\neg\varphi\to(\varphi\to\psi)\);
  \item \((\neg\psi\to\neg\varphi)\to(\varphi\to\psi)\);
  \item \((\varphi\to\psi)\to(\neg\psi\to\neg\varphi)\);
  \item \((\varphi\to\psi)\to((\neg\varphi\to\psi)\to\psi)\).
\end{enumerate}
\end{hblemma}

\begin{hblemma}{thm:p02-w02b-p02-010-rule-images}{proved=annex:p02-w02b-p02-010-rule-images}{Hilbert images of the natural-deduction rules}
For all \(A,B,C\in\mathsf{Fm}_{\{\neg,\to\}}\), writing \(\vdash\) for
\(\derives_{\mathsf H_{\mathsf{CPL}}}\):
\[
\begin{array}{ll}
 \text{(i)}\ \{A,B\}\vdash\neg(A\to\neg B);
 &\text{(ii)}\ \{\neg(A\to\neg B)\}\vdash A\ \text{and}\
   \{\neg(A\to\neg B)\}\vdash B;\\[2pt]
 \text{(iii)}\ \{A\}\vdash\neg A\to B\ \text{and}\ \{B\}\vdash\neg A\to B;
 &\text{(iv)}\ \{\neg A\to B,\ A\to C,\ B\to C\}\vdash C;\\[2pt]
 \text{(v)}\ \{A\to\bot_{\mathsf H}\}\vdash\neg A;
 &\text{(vi)}\ \{A,\neg A\}\vdash\bot_{\mathsf H};\\[2pt]
 \text{(vii)}\ \{\bot_{\mathsf H}\}\vdash C;
 &\text{(viii)}\ \{\neg A\to\bot_{\mathsf H}\}\vdash A.
\end{array}
\]
\end{hblemma}

Items (i)--(viii) are the translated forms of conjunction introduction and
elimination, disjunction introduction and elimination, negation
introduction and elimination, falsum elimination and RAA: with
\(A=\operatorname{tr}_{\mathsf H}(\alpha)\) and
\(B=\operatorname{tr}_{\mathsf H}(\beta)\), the formula
\(\neg(A\to\neg B)\) is \(\operatorname{tr}_{\mathsf H}(\alpha\land\beta)\)
and \(\neg A\to B\) is \(\operatorname{tr}_{\mathsf H}(\alpha\lor\beta)\).

\subsection{Natural-deduction facts used by the maps}
\label{sec:p02-w02b-p02-010-nd}

\begin{hblemma}{thm:p02-w02b-p02-010-nd-structural}{proved=annex:p02-w02b-p02-010-nd-structural}{Padding, unrestricted discharge and composition in natural deduction}
Let \(\Gamma,\Delta,\Sigma\subseteq\mathsf{Fm}_{\mathsf{CPL}}\) and
\(\alpha,\beta,\chi\in\mathsf{Fm}_{\mathsf{CPL}}\).
\begin{enumerate}
  \item (Padding.) For every derivation \(D\) with root \(\chi\) there is a
    derivation \(D'\) with root \(\chi\), containing an open occurrence of
    \(\alpha\), and with
    \(\operatorname{oa}(D')=\operatorname{oa}(D)\cup\{\alpha\}\).
  \item (Discharge as a closure property.)  If
    \(\Gamma\cup\{\alpha\}\derives_{\mathsf{ND}_{\mathsf{CPL}}}\beta\) then
    \(\Gamma\derives_{\mathsf{ND}_{\mathsf{CPL}}}\alpha\to\beta\); if
    \(\Gamma\cup\{\alpha\}\derives_{\mathsf{ND}_{\mathsf{CPL}}}\bottom\) then
    \(\Gamma\derives_{\mathsf{ND}_{\mathsf{CPL}}}\neg\alpha\); if
    \(\Gamma\cup\{\neg\alpha\}\derives_{\mathsf{ND}_{\mathsf{CPL}}}\bottom\)
    then \(\Gamma\derives_{\mathsf{ND}_{\mathsf{CPL}}}\alpha\); and if
    \(\Gamma\derives_{\mathsf{ND}_{\mathsf{CPL}}}\alpha\lor\beta\),
    \(\Delta\cup\{\alpha\}\derives_{\mathsf{ND}_{\mathsf{CPL}}}\chi\) and
    \(\Sigma\cup\{\beta\}\derives_{\mathsf{ND}_{\mathsf{CPL}}}\chi\), then
    \(\Gamma\cup\Delta\cup\Sigma\derives_{\mathsf{ND}_{\mathsf{CPL}}}\chi\).
    This holds whether or not the witnessing derivation of the premise uses
    the assumption.
  \item (Composition.) If \(\Delta\derives_{\mathsf{ND}_{\mathsf{CPL}}}\sigma\)
    for every \(\sigma\in\Sigma\) and
    \(\Gamma\cup\Sigma\derives_{\mathsf{ND}_{\mathsf{CPL}}}\chi\), then
    \(\Gamma\cup\Delta\derives_{\mathsf{ND}_{\mathsf{CPL}}}\chi\).
\end{enumerate}
\end{hblemma}

\begin{hblemma}{thm:p02-w02b-p02-010-nd-translation}{proved=annex:p02-w02b-p02-010-nd-translation}{Each formula is natural-deduction equivalent to its Hilbert translation}
For every \(\alpha\in\mathsf{Fm}_{\mathsf{CPL}}\),
\[
  \{\alpha\}\derives_{\mathsf{ND}_{\mathsf{CPL}}}\operatorname{tr}_{\mathsf H}(\alpha)
  \qquad\text{and}\qquad
  \{\operatorname{tr}_{\mathsf H}(\alpha)\}\derives_{\mathsf{ND}_{\mathsf{CPL}}}\alpha .
\]
\end{hblemma}

\subsection{The four maps}
\label{sec:p02-w02b-p02-010-maps}

\begin{hbtheorem}{thm:p02-w02b-p02-010-hilbert-to-nd}{proved=annex:p02-w02b-p02-010-hilbert-to-nd}{From Hilbert derivations to natural deduction}
Let \(\Gamma\subseteq\mathsf{Fm}_{\{\neg,\to\}}\) and let
\(\theta_0,\ldots,\theta_n\) be a Hilbert derivation from \(\Gamma\).  There
are \(\mathsf{ND}_{\mathsf{CPL}}\) derivations \(D_0,\ldots,D_n\), constructed
line by line, such that \(D_i\) has root \(\theta_i\) and
\(\operatorname{oa}(D_i)\subseteq\Gamma\).  A premise line becomes a one-node
assumption, an axiom line becomes a fixed closed derivation of that
instance, and an MP line becomes \(\to E\) applied to relabelled copies of
the two earlier trees.  Consequently
\(\Gamma\derives_{\mathsf H_{\mathsf{CPL}}}\theta\) implies
\(\Gamma\derives_{\mathsf{ND}_{\mathsf{CPL}}}\theta\).  The derivations
produced use only the rules assumption, \(\land I\), \(\land E_1\),
\(\land E_2\), \(\to I\), \(\to E\), \(\neg E\) and RAA.
\end{hbtheorem}

\begin{hbtheorem}{thm:p02-w02b-p02-010-nd-to-hilbert}{proved=annex:p02-w02b-p02-010-nd-to-hilbert}{From natural deduction to Hilbert derivations}
For every \(\mathsf{ND}_{\mathsf{CPL}}\) derivation \(D\) with root
\(\alpha\) there is a Hilbert derivation, constructed by recursion on \(D\),
of \(\operatorname{tr}_{\mathsf H}(\alpha)\) from
\(\operatorname{tr}_{\mathsf H}[\operatorname{oa}(D)]\).  Each rule
application of \(D\) is replaced by the matching item of
\cref{thm:p02-w02b-p02-010-rule-images}, \(\to E\) by MP, and \(\to I\) by
the deduction theorem.  Consequently, for every
\(\Gamma\subseteq\mathsf{Fm}_{\mathsf{CPL}}\),
\(\Gamma\derives_{\mathsf{ND}_{\mathsf{CPL}}}\alpha\) implies
\(\operatorname{tr}_{\mathsf H}[\Gamma]\derives_{\mathsf H_{\mathsf{CPL}}}
\operatorname{tr}_{\mathsf H}(\alpha)\).
\end{hbtheorem}

\begin{hbtheorem}{thm:p02-w02b-p02-010-hilbert-to-lk}{proved=annex:p02-w02b-p02-010-hilbert-to-lk}{From Hilbert derivations to sequent derivations}
Let \(\Gamma\subseteq\mathsf{Fm}_{\{\neg,\to\}}\) and let
\(\theta_0,\ldots,\theta_n\) be a Hilbert derivation from \(\Gamma\).  For
each \(i\) there is a finite multiset \(\Pi_i\) with
\(\operatorname{supp}(\Pi_i)\subseteq\Gamma\) and an
\(\mathsf{LK}_{\mathsf{CPL}}\) derivation of \(\Pi_i\Rightarrow\theta_i\),
constructed line by line: a premise line is an identity sequent, an axiom
line is a fixed cut-free derivation of \(\Rightarrow\theta_i\), and an MP line
is two cuts against the sequent \(\varphi,\varphi\to\psi\Rightarrow\psi\).
Consequently, for every finite \(\Gamma_0\subseteq\mathsf{Fm}_{\{\neg,\to\}}\)
with \(\operatorname{supp}(\Pi_n)\subseteq\Gamma_0\), the sequent
\(\Gamma_0\Rightarrow\theta_n\) is derivable in \(\mathsf{LK}_{\mathsf{CPL}}\);
in particular, for finite \(\Gamma\),
\(\Gamma\derives_{\mathsf H_{\mathsf{CPL}}}\theta\) implies that
\(\Gamma\Rightarrow\theta\) is derivable.
\end{hbtheorem}

\begin{hbtheorem}{thm:p02-w02b-p02-010-lk-to-hilbert}{proved=annex:p02-w02b-p02-010-lk-to-hilbert}{From sequent derivations to Hilbert derivations}
Let \(\Gamma\Rightarrow\Delta\) be derivable in \(\mathsf{LK}_{\mathsf{CPL}}\),
and let \(D_{\Delta}=\bigvee\Delta\) be the disjunction of
\cref{def:p07-sequent-nd-translation}, with \(D_{\varnothing}=\bottom\).
Then
\(\operatorname{tr}_{\mathsf H}[\operatorname{supp}(\Gamma)]
\derives_{\mathsf H_{\mathsf{CPL}}}\operatorname{tr}_{\mathsf H}(D_{\Delta})\).
The Hilbert derivation is obtained by applying the sequent-to-natural-deduction
construction of \cref{annex:p7.calculi.3} and then the construction of
\cref{thm:p02-w02b-p02-010-nd-to-hilbert}; for a single succedent formula
\(\alpha\) the conclusion reads
\(\operatorname{tr}_{\mathsf H}[\operatorname{supp}(\Gamma)]
\derives_{\mathsf H_{\mathsf{CPL}}}\operatorname{tr}_{\mathsf H}(\alpha)\).
\end{hbtheorem}

\begin{hbcorollary}{thm:p02-w02b-p02-010-hilbert-equivalence}{proved=annex:p02-w02b-p02-010-hilbert-equivalence}{Hilbert legs of the derivability equivalence}
For every \(\Gamma\subseteq\mathsf{Fm}_{\mathsf{CPL}}\) and
\(\alpha\in\mathsf{Fm}_{\mathsf{CPL}}\):
\begin{enumerate}
  \item \(\Gamma\derives_{\mathsf{ND}_{\mathsf{CPL}}}\alpha\) iff
    \(\operatorname{tr}_{\mathsf H}[\Gamma]\derives_{\mathsf H_{\mathsf{CPL}}}
    \operatorname{tr}_{\mathsf H}(\alpha)\);
  \item if \(\Gamma\) is finite, the sequent \(\Gamma\Rightarrow\alpha\) is
    derivable in \(\mathsf{LK}_{\mathsf{CPL}}\) iff
    \(\operatorname{tr}_{\mathsf H}[\Gamma]\derives_{\mathsf H_{\mathsf{CPL}}}
    \operatorname{tr}_{\mathsf H}(\alpha)\);
  \item if \(\Gamma\cup\{\alpha\}\subseteq\mathsf{Fm}_{\{\neg,\to\}}\), then
    \(\Gamma\derives_{\mathsf H_{\mathsf{CPL}}}\alpha\) iff
    \(\Gamma\derives_{\mathsf{ND}_{\mathsf{CPL}}}\alpha\), and, for finite
    \(\Gamma\), iff \(\Gamma\Rightarrow\alpha\) is derivable in
    \(\mathsf{LK}_{\mathsf{CPL}}\).
\end{enumerate}
\end{hbcorollary}

Together with \cref{thm:prop-nd-soundness,thm:prop-completeness} for the
semantic leg and with \cref{thm:p07-lk-nd-equivalence} restricted to
propositional derivations, this corollary yields every equivalence displayed
in \cref{thm:p02-calculi-derivability-equivalence}.  The corollary is proved
syntactically: no step uses soundness or completeness of any calculus.

\begin{hbmethod}{def:p02-w02b-p02-010-procedure}{Translating a derivation between the three calculi}
\emph{Input:} a derivation in one of \(\mathsf H_{\mathsf{CPL}}\),
\(\mathsf{ND}_{\mathsf{CPL}}\), \(\mathsf{LK}_{\mathsf{CPL}}\), and a target
calculus.
\begin{enumerate}
  \item \(\mathsf H\to\mathsf{ND}\): process lines in order; for line
    \(\theta_i\) take an assumption leaf, the fixed closed derivation of the
    axiom instance (with the A1 case split \(\varphi\neq\psi\) or
    \(\varphi=\psi\)), or \(\to E\) beneath copies of the trees for its two
    MP premises, relabelled apart.
  \item \(\mathsf{ND}\to\mathsf H\): recurse from the leaves to the root of
    the tree; at a rule application, concatenate the Hilbert derivations of
    the translated premises and append the derivation of the matching item
    of \cref{thm:p02-w02b-p02-010-rule-images}; for a discharging rule, first
    apply the deduction theorem to the discharged translated assumption.
  \item \(\mathsf H\to\mathsf{LK}\): process lines in order as in step 1,
    using identity sequents, the three fixed cut-free axiom derivations, and
    two cuts for MP; finish with left contraction and left weakening.
  \item \(\mathsf{LK}\to\mathsf H\): apply the sequent-to-natural-deduction
    construction of \cref{annex:p7.calculi.3}, then step 2.
  \item \(\mathsf{ND}\to\mathsf{LK}\) and \(\mathsf{LK}\to\mathsf{ND}\): the
    two constructions of \cref{annex:p7.calculi.3}.
\end{enumerate}
\emph{Side conditions:} every discharge in a produced natural-deduction tree
must be nonvacuous; the constructions guarantee this by the padding detour
of \cref{thm:p02-w02b-p02-010-nd-structural}(1) and the A1 case split.
Discharge labels in copied subtrees must be renamed apart before grafting or
\(\to E\).  The Hilbert side has no discharge, so no side condition arises
there beyond the formula shapes of the schemas and of MP.
\emph{Output:} a derivation in the target calculus of the same formula or,
for maps into \(\mathsf H_{\mathsf{CPL}}\), of its
\(\operatorname{tr}_{\mathsf H}\)-image, from translated assumptions.
\end{hbmethod}

\subsection{Worked example}
\label{sec:p02-w02b-p02-010-example}

\begin{hbexample}{def:p02-w02b-p02-010-example}{The conjunction projection through all three calculi}
Take the natural-deduction derivation of \((P\land Q)\to P\) in
\cref{fig:p02-nd-lk-correspondence}: an assumption \([P\land Q]^1\),
\(\land E_1\) to \(P\), and \(\to I\) discharging label \(1\).  Its
translation is
\(\operatorname{tr}_{\mathsf H}((P\land Q)\to P)=\neg(P\to\neg Q)\to P\).

\emph{Natural deduction to Hilbert.}  The leaf becomes the one-line
derivation of \(\neg(P\to\neg Q)\) from itself.  The \(\land E_1\) step
appends the derivation of item (ii) of
\cref{thm:p02-w02b-p02-010-rule-images}, whose lines are
\[
\begin{array}{rll}
1.&\neg(P\to\neg Q)&\text{premise}\\
2.&\neg P\to(P\to\neg Q)&\text{\cref{thm:p02-w02b-p02-010-hilbert-theorems}(5)}\\
3.&(\neg P\to(P\to\neg Q))\to(\neg(P\to\neg Q)\to\neg\neg P)&\text{\cref{thm:p02-w02b-p02-010-hilbert-theorems}(7)}\\
4.&\neg(P\to\neg Q)\to\neg\neg P&\text{MP }2,3\\
5.&\neg\neg P&\text{MP }1,4\\
6.&\neg\neg P\to P&\text{\cref{thm:p02-w02b-p02-010-hilbert-theorems}(3)}\\
7.&P&\text{MP }5,6
\end{array}
\]
where lines 2, 3 and 6 stand for their explicit empty-premise derivations
from the annex, inserted in place.  The \(\to I\) step applies the deduction
theorem to this derivation, discharging the premise
\(\neg(P\to\neg Q)\), and yields
\(\derives_{\mathsf H_{\mathsf{CPL}}}\neg(P\to\neg Q)\to P\).

\emph{Hilbert back to natural deduction.}  Applied to that Hilbert
derivation, \cref{thm:p02-w02b-p02-010-hilbert-to-nd} gives a closed
derivation of \(\neg(P\to\neg Q)\to P\).  To return to the original
formula, the implication case of
\cref{thm:p02-w02b-p02-010-nd-translation} is used: assume
\([P\land Q]^2\), derive \(\neg(P\to\neg Q)\) from it by the conjunction case
of the same lemma, apply \(\to E\) with the closed derivation to obtain
\(P\), and discharge label \(2\).  The result is a closed derivation of
\((P\land Q)\to P\).  It is a different tree from the one in the figure: the
maps preserve the conclusion, not the derivation.

\emph{Hilbert to sequents.}  Applied to the same Hilbert derivation,
\cref{thm:p02-w02b-p02-010-hilbert-to-lk} produces an
\(\mathsf{LK}_{\mathsf{CPL}}\) derivation of
\(\Rightarrow\neg(P\to\neg Q)\to P\) that uses Cut at every MP line, whereas
the derivation of \(\Rightarrow(P\land Q)\to P\) in the figure is cut-free.
\end{hbexample}

\subsection{Limitations and a failure case}
\label{sec:p02-w02b-p02-010-limits}

\begin{hbwarning}{def:p02-w02b-p02-010-vacuous}{Vacuous discharge defeats the obvious translation of A1}
The usual natural-deduction proof of the A1 instance
\(\varphi\to(\psi\to\varphi)\) assumes \(\varphi\), discharges a never-used
\(\psi\), and then discharges \(\varphi\).  On
\texttt{sheet:\allowbreak nd-cpl} that is not a derivation: a
discharge needs an open occurrence of its formula.  The padding detour
\(\varphi,\psi\vdash\varphi\land\psi\vdash\varphi\) supplies the occurrence
of \(\psi\) when \(\psi\neq\varphi\).  It fails for the instance
\(P\to(P\to P)\), because the inner \(\to I\) discharges every open
occurrence of \(P\), leaving none for the outer one.  The construction
therefore treats \(\varphi=\psi\) separately: it conjoins the assumption
\(P\) with the closed derivation of \(P\to P\), projects to \(P\to P\), and
discharges \(P\) once.  A translation that ignores the side condition yields
a tree that is not a derivation.
\end{hbwarning}

\begin{hbwarning}{def:p02-w02b-p02-010-size}{The maps preserve conclusions, not size or proof shape}
None of the constructions preserves derivation size up to a constant
factor.  A Hilbert line used by several later MP steps is copied into each
natural-deduction or sequent subtree that needs it: at an MP line the new
tree contains complete copies of both premise trees.  Tree sizes therefore
add along every MP step, and for a sequence of \(n\) lines the construction
guarantees no bound better than one exponential in \(n\).  In
the other direction each application of the deduction theorem replaces
every line of the derivation it acts on by a segment of up to five lines, so
nested discharges multiply length.  The sequent derivations produced from
Hilbert derivations contain Cut at every MP line; a cut-free derivation of
the same sequent exists by cut elimination, which is a separate Part~7
result and is not used here.  The translation \(\operatorname{tr}_{\mathsf H}\)
also depends on the fixed variable \(p_0\): it preserves each formula's truth
function, but \(\operatorname{tr}_{\mathsf H}(\bottom)\) is not \(\bottom\),
and the ND--Hilbert round trip returns a derivation of a formula only
up to the natural-deduction equivalence of
\cref{thm:p02-w02b-p02-010-nd-translation}.
\end{hbwarning}

\subsection{Relationships and sources}
\label{sec:p02-w02b-p02-010-relations}

The three-way comparison of the calculi by derivation shape and use is
\cref{tab:p02-proof-systems}, and one natural-deduction/sequent
correspondence is drawn in \cref{fig:p02-nd-lk-correspondence}.  The maps of
this section are the propositional counterparts of the first-order proof
bridges \cref{thm:p02-hfa-proof-bridges}; that construction obtains its
propositional kernels from \cref{thm:p02-calculi-derivability-equivalence},
whose Hilbert legs are now given by
\cref{thm:p02-w02b-p02-010-hilbert-equivalence} without semantic
completeness.  The rejection calculus of
\cref{sec:p02-classical-rejection} takes the Hilbert derivations of
\cref{def:p02-hilbert-derivation} as its assertion judgment, and its
classical adequacy is the empty-premise case of these equivalences.

Gentzen proved the equivalence of his natural-deduction calculus
\(\mathsf{NK}\), his sequent calculus \(\mathsf{LK}\) and a calculus
\(\mathsf{LHK}\) modelled on Hilbert's formalism by the cycle
\(\mathsf{LHK}\to\mathsf{NK}\to\mathsf{LK}\to\mathsf{LHK}\)
\cite[Abschnitt~V, \S 6, p.~429]{Gentzen1935b}, counting two derivations
equivalent when their end formulas or end sequents are equivalent in the
sense fixed in \cite[Abschnitt~V, \S 1, pp.~417--418]{Gentzen1935b}.
Kleene proves the deduction theorem for his propositional calculus by
course-of-values induction on the length of the given deduction
\cite[\S 21, Theorem~1, p.~90]{Kleene1952}, and translates between his
Hilbert-style system and Gentzen's sequent system \(\mathsf{G1}\) in both
directions \cite[\S 77, Theorems~46 and 47, pp.~445--446]{Kleene1952}; his
translation of an empty succedent uses a negated self-implication, as
\(\operatorname{tr}_{\mathsf H}(\bottom)\) does here.  Neither source uses
the calculi of this book; for instance, Kleene's Axiom Schema 1b, displayed
at \cite[\S 21, p.~92]{Kleene1952}, is A2 with its two antecedents
interchanged.  The constructions and proofs in this section are written
directly for \texttt{sheet:\allowbreak hilbert-cpl},
\texttt{sheet:\allowbreak nd-cpl} and the propositional rows
of \texttt{sheet:\allowbreak lk}; the axiom schemas follow the
attribution on the sheet \cite{Mendelson1997}.
